\begin{proof}
(i) By Theorem \ref{thm:c11-classes} (iii) the two-body part is $\Upsilon_\xi$ on the single pair $(A|_{W_{\mathbb R}^\perp},A|_{W_{\mathbb R}^\perp})$ of $O$, whose leg sets are those of Table \ref{tab:c11-defectpairs}: $\lambda|_{S_-\otimes W^*}$, $A|_V\,\lambda|_{S_-\otimes W^*}$, $A|_{W_{\mathbb R}^\perp}\lambda|_{S_+\otimes W}$ and $c\,\lambda^*|_{S_+\otimes W^*}$ for $\xi\in S_+\otimes W$. The single output $Y$ has the colour factor $(1\otimes\mathrm{ad}_Y)(-,-)^{-1}$, which the colour pairing of $O$ contracts to $\operatorname{tr}_{\mathfrak g}\mathrm{ad}_Y=0$. The other three leg sets have $D=0$, so they carry no poles and no finite special vertex (Proposition \ref{prop:c11-oddbound}), and they give $\mathcal B_1$, $\mathcal B_2$ (outputs $A|_{W_{\mathbb R}^\perp}\lambda$), $\mathcal B_3$ (outputs $A|_V\lambda$) and $\mathcal B_4$ (outputs $c\,\lambda^*$). Their chains have three edges and, at $D=0$, constant integrands, so each chain gives a rational multiple of $\tfrac12\log\tfrac32$ (Theorem \ref{thm:c11-oddsimplex} (iv)); the exact evaluation of the chains in all placements, in both ways of Remark \ref{rem:c11-oddmethod}, gives \eqref{eq:c11-Q1O}, with the same coefficient for $\mathcal B_2$ and $\mathcal B_3$. (ii) By Corollary \ref{cor:c7-Q} the functionals at $(w_1,w_2)=0$ are the multiples of $\mathcal B_1$; the one-body term is $-t_\perp\mathcal B_1$ units (proof of Lemma \ref{lem:c11-cocycle} (i)); the changes are Propositions \ref{prop:c11-freexi} and \ref{prop:c11-freedom}. (iii) $L_1^{(0)}$ is given by the arity-zero two-body kernels, the cycles with one special edge of the proof of Theorem \ref{thm:c10-brst}, of the pairs $(A|_{W_{\mathbb R}^\perp},A|_{W_{\mathbb R}^\perp})$ of $O$ and $(A|_{W_{\mathbb R}^\perp},\lambda)$ of $\mathcal B_1$, evaluated in the same way; then Lemma \ref{lem:c11-cocycle} (ii). (v) By (i) and (ii), $T_2$ being the only term of the solution of Proposition \ref{prop:c11-oddsol} (i) with the leg $A^*|_{W_{\mathbb R}^\perp}$. (iv) With $s\mathcal B_3=-(w_1,w_2)\int\omega\,(\rho(dc)\xi,\lambda)$ in this sector, the three terms are $4L\,s\mathcal B_3$, $2\,s\mathcal B_3$ and $2\kappa_O\,s\mathcal B_3=-(2+4L)\,s\mathcal B_3$ units (proof of Lemma \ref{lem:c11-cocycle} (ii)).
\end{proof}
